\chapter{How NatQuery Works}
\label{ch:how}

This chapter follows a question through NatQuery and explains each mechanism: the request path, the
policy and role view, the SQL enforcer, the repair loop, schema selection for large databases,
multi-schema support with row-level security, execution limits, the audit log and authentication.

\section{The Request Path}

Table~\ref{tab:path} lists the thirteen steps of Figure~\ref{fig:arch} with what happens at each and
what happens when it fails. A failure at any step ends the request with a generic message and an opaque
reference id; the details go only to the audit log.

\begin{table}[htbp]
\centering
\caption{The thirteen steps of a request.}
\label{tab:path}
\small
\renewcommand{\arraystretch}{1.15}
\begin{tabularx}{\textwidth}{L{0.7cm} L{3.1cm} Y L{3.3cm}}
\toprule
\hdr \textbf{\#} & \textbf{Step} & \textbf{What happens} & \textbf{On failure} \\
\midrule
1 & Session check & The token is validated and the identity re-read from the provider, so a disabled account or a changed role takes effect at once & Request refused \\
2 & Role selection & The highest-priority role the user holds, or the one the user selected; never a union of roles & Refused if no policy role \\
3 & Input guard, limits & Length and character checks, injection patterns, per-user rate and model-call limits & Refused, audited \\
4 & Role view & Tables and columns the role may use, masks and row filters, built from the policy & Refused if the view is empty \\
5 & Schema selection & Small views go whole; large views are ranked and trimmed (Section~\ref{sec:select}) & --- \\
6 & Prompt & Output contract, access rules, schema with value hints, and the question & --- \\
7 & Model call & One SQL statement is requested; the reply is cleaned (fences, comments) & Provider error: generic message \\
8 & Enforcement & Parse, check, rewrite, regenerate (Section~\ref{sec:enforcer}) & Fixable: repair loop; otherwise refused \\
9 & Audit record & ``authorized'' record with model SQL, executed SQL, tables, masks, filters & \textbf{Nothing executes} \\
10 & Execution & Read-only, time- and row-bounded, under the database role and with identity settings & Sanitised error to repair loop \\
11 & Result guard & Pattern redaction and size caps on values & --- \\
12 & Audit record & ``completed'' record with outcome and timing & --- \\
13 & Answer & Rows, columns, notes (for example masking or truncation) and the reference id & --- \\
\bottomrule
\end{tabularx}
\end{table}

\section{Policy and Role View}

The policy is one YAML file. A role is a block in it; the names ``student'', ``faculty'' or ``admin''
mean nothing to the code. Table~\ref{tab:policykeys} lists the keys, and Listing~\ref{lst:policy} shows an
excerpt of the ERP policy.

\begin{table}[htbp]
\centering
\caption{Policy keys.}
\label{tab:policykeys}
\small
\renewcommand{\arraystretch}{1.15}
\begin{tabularx}{\textwidth}{L{5.3cm} Y}
\toprule
\hdr \textbf{Key} & \textbf{Meaning} \\
\midrule
\code{roles.<name>.inherits} & Roles whose grants this role starts from; child blocks override per table \\
\code{roles.<name>.priority} & Decides which role a multi-role user runs under \\
\code{tables.<t>.columns} & Visible columns, or \code{"*"}; \code{schema.*} grants a whole schema \\
\code{tables.<t>.deny_columns} & Columns removed from the above \\
\code{tables.<t>.masked} & \code{column: strategy}; value replaced at SELECT time, selectable but never filterable \\
\code{tables.<t>.row_filter} & Condition ANDed onto every access to the table; \code{:placeholders} bound from the identity \\
\code{global_deny} & Tables, \code{table.column} and \code{*.column} patterns no role can see, admin included \\
\code{allowed_functions} & Extra SQL functions beyond the built-in allow-list \\
\code{defaults} / per-role limits & \code{max_rows} (500), \code{statement_timeout_ms} (5000), \code{max_joins} (6), subquery depth, SQL length \\
\code{permissions} & \code{can_view_sql}, \code{can_view_plan}, \code{can_manage_users}, \code{can_view_audit}, \code{can_refresh_schema} \\
\code{db_role}, \code{db_session_vars} & Database role to drop to, and identity attributes to hand to row-level security \\
\bottomrule
\end{tabularx}
\end{table}

\begin{lstlisting}[caption={Excerpt of the ERP policy (student role). Row filters use placeholders bound from the user's identity.}, label={lst:policy}]
global_deny:
  tables: [auth.session, auth.otp]
  columns: ['*.password_hash', '*.token_hash', people.person_identity.aadhaar_no]
db_session_vars:
  erp.role: '@role'
  erp.student_id: student_id
roles:
  student:
    inherits: [erp_base]
    priority: 10
    tables:
      people.student:
        columns: '*'
        row_filter: student_id = :student_id
      acad.result:
        columns: '*'
        row_filter: registration_id IN (SELECT registration_id
                    FROM acad.registration WHERE student_id = :student_id)
\end{lstlisting}

The role view is built from the policy and the extracted schema. It contains only granted tables and
columns, the masks and the row filters, and a prompt-facing copy for the model. Because the view is the
only schema the model and the enforcer ever see, a hidden table is not forbidden: it does not exist.

\section{The SQL Enforcer}
\label{sec:enforcer}

The enforcer turns the model's text into SQL that is safe to run. It works on a syntax tree built with the
\texttt{sqlglot} parser \citep{sqlglot}, in the stages of Figure~\ref{fig:enforcer}, and finally
\emph{regenerates} the SQL from the tree, so that no byte the model wrote reaches the database.

\begin{figure}[htbp]
\centering
\resizebox{0.95\textwidth}{!}{%
\begin{tikzpicture}[
  st/.style={draw=nqblue, fill=nqlight, rounded corners=2pt, minimum width=3.1cm, minimum height=1.2cm, align=center, font=\footnotesize, text width=2.9cm},
  fl/.style={-{Latex[length=2mm]}, thick, nqblue}]
\node[draw=nqred, dashed, rounded corners=2pt, minimum width=3.1cm, minimum height=1.2cm, align=center, font=\footnotesize] (in) at (0,0) {Model text (untrusted)};
\node[st] (a) at (4.1,0) {\textbf{1 Parse}\\one SELECT only};
\node[st] (b) at (8.2,0) {\textbf{2 Allow-list}\\constructs and functions};
\node[st] (c) at (12.3,0) {\textbf{3 Resolve}\\names against the role view};
\node[st] (d) at (12.3,-2) {\textbf{4 Row filters}\\wrap every table reference};
\node[st] (e) at (8.2,-2) {\textbf{5 Masks}\\display-only columns};
\node[st] (f) at (4.1,-2) {\textbf{6 Limits}\\LIMIT, joins, depth};
\node[draw=nqgreen, rounded corners=2pt, minimum width=3.1cm, minimum height=1.2cm, align=center, font=\footnotesize] (out) at (0,-2) {\textbf{7 Regenerate}\\safe SQL for the database};
\foreach \x/\y in {in/a,a/b,b/c,c/d,d/e,e/f,f/out} \draw[fl] (\x) -- (\y);
\end{tikzpicture}}
\caption{Stages of the SQL enforcer. A rule violation at any stage ends in a refusal or, if fixable, a repair request.}
\label{fig:enforcer}
\end{figure}

\begin{table}[htbp]
\centering
\caption{What the enforcer rejects or changes.}
\label{tab:enforcer}
\small
\renewcommand{\arraystretch}{1.15}
\begin{tabularx}{\textwidth}{L{3.1cm} Y L{3.9cm}}
\toprule
\hdr \textbf{Rule} & \textbf{Effect} & \textbf{Example blocked} \\
\midrule
Statement type & Only SELECT and set operations; no INSERT, UPDATE, DELETE, DDL, INTO or row locks; one statement only & \code{DROP TABLE x}; stacked statements \\
Constructs & No recursive CTEs, LATERAL, table functions, bind parameters or catalog casts & \code{WITH RECURSIVE ...} \\
Functions & Allow-list only (aggregates, \code{COALESCE}, \code{CASE}, date and string functions) & \code{pg_sleep}, \code{pg_read_file}, \code{lo_import} \\
Name resolution & Every table and column must exist in the role view; hidden and non-existent names give the same error & Hidden table, hidden column \\
Whole-row references & A table name used as a value is refused, since it would return hidden columns & \code{SELECT employee::text FROM employee} \\
Row filters & Each table reference, in joins, subqueries, CTEs and UNIONs, is replaced by a filtered subquery; identity values are validated and bound as literals & \code{OR TRUE}, self-join, UNION bypasses \\
Masked columns & May be selected and are replaced by the mask; any use in WHERE, JOIN, GROUP BY, ORDER BY, functions or aggregates is refused (no oracle) & \code{WHERE} or \code{ORDER BY} on a masked column \\
Resource limits & No cross joins (including \code{ON true}); joins, nesting, SQL length and syntax-tree size capped; \code{LIMIT} enforced & Cartesian products, \code{repeat('a', 10^9)} \\
\bottomrule
\end{tabularx}
\end{table}

\section{A Worked Example}
\label{sec:example}

Take the question \emph{``Which library fines have I not paid yet, and for which books?''} asked by a
student of the synthetic ERP. The model returns an ordinary query that knows nothing of the user's
rights. The enforcer replaces the two personal tables by filtered subqueries bound to the caller's
\code{person_id}, adds a row cap, and leaves the shared catalogue tables untouched
(Listing~\ref{lst:example}). The student cannot see another person's fines even if the model forgets to
filter.

\begin{lstlisting}[caption={Model SQL and executed SQL for one question (identity value from the synthetic data).}, label={lst:example}]
-- SQL written by the model
SELECT d.reason, d.amount, b.title
FROM lib.library_due d
JOIN lib.loan l ON l.loan_id = d.loan_id
JOIN lib.book_copy bc ON bc.copy_id = l.copy_id
JOIN lib.book b ON b.book_id = bc.book_id
WHERE d.status = 'UNPAID'

-- SQL regenerated by the enforcer and executed
SELECT d.reason AS reason, d.amount AS amount, b.title AS title
FROM (SELECT * FROM lib.library_due WHERE person_id = '1487') AS d
JOIN (SELECT * FROM lib.loan WHERE person_id = '1487') AS l ON l.loan_id = d.loan_id
JOIN lib.book_copy AS bc ON bc.copy_id = l.copy_id
JOIN lib.book AS b ON b.book_id = bc.book_id
WHERE d.status = 'UNPAID'
LIMIT 200
\end{lstlisting}

\section{The Repair Loop}

Models often make small mistakes that a precise message can fix. When the enforcer or the database
rejects a statement, the error is classified (Table~\ref{tab:errors}), stripped of anything sensitive and
sent back with the failed SQL, and the model tries again. At most two repairs are made (configurable); a
statement that still fails is refused. Policy violations are never ``repaired'' into a way around the
policy: the hint tells the model to use only what the schema shows.

\begin{table}[htbp]
\centering
\caption{Error classes used to give the model a targeted hint.}
\label{tab:errors}
\small
\renewcommand{\arraystretch}{1.15}
\begin{tabularx}{\textwidth}{L{3.8cm} Y}
\toprule
\hdr \textbf{Class} & \textbf{Typical cause and hint} \\
\midrule
\code{SYNTAX_ERROR} & Malformed SQL; return one valid statement \\
\code{UNDEFINED_TABLE} / \code{UNDEFINED_COLUMN} & Name not in the schema; use only listed names, or answer ``unanswerable'' \\
\code{AMBIGUOUS_COLUMN} & Column exists in several joined tables; qualify it \\
\code{GROUP_BY_ERROR} & Selected column neither grouped nor aggregated \\
\code{JOIN_ERROR} & Missing or invalid join condition; every JOIN needs an ON clause \\
\code{TYPE_ERROR} & Operator or cast on the wrong type \\
\code{POLICY} & Rule broken (cross join, display-only column misuse, unsupported construct) \\
\code{UNKNOWN} & Anything else; generic retry \\
\bottomrule
\end{tabularx}
\end{table}

\section{Schema Selection for Large Databases}
\label{sec:select}

A university ERP has far more schema than fits in a free-tier prompt. NatQuery therefore sends a small
view whole (up to 250 columns and 20 tables) and, for larger ones, selects the relevant tables in four
steps.

\begin{enumerate}
  \item \textbf{Rank.} Each table gets a score from the words it shares with the question,
        weighted by rarity so that words like \emph{id} or \emph{date} count for little:
        \[
          \mathrm{score}(t) = \sum_{w \in q \cap T_t} \omega_{t,w}\,\log\!\Bigl(1 + \tfrac{N}{\mathrm{df}_w}\Bigr),
        \]
        where $T_t$ are the words of table $t$ (name, description, synonyms, columns), $N$ the number
        of tables, $\mathrm{df}_w$ the number of tables containing $w$, and $\omega_{t,w}$ is 3 for a
        word in the table name, 2 for the description or synonyms and 1 for a column name.
  \item \textbf{Connect.} Take the top five tables, add their foreign-key neighbours and the tables
        that bridge them (up to three hops). If more than 20 result, keep the best ranked, then those
        connected to the most chosen tables.
  \item \textbf{Add lookups.} Small tables (at most eight columns) referenced by a chosen table, such as
        department, program or session, are added: they are cheap and are how an id becomes a name.
  \item \textbf{Hint values.} Columns with few distinct values show them (for example
        \code{status -- values: 'FAILED', 'SUCCESS'}), so the model does not guess spellings. Hints are
        collected once by an administrator, shown only for columns the role may see, and never for
        masked or secret-looking columns.
\end{enumerate}

If the model still answers ``unanswerable'' on a trimmed schema, the request is retried once with a wider
selection (top ten tables, up to 32) before the user is told the data is not available.

\section{Several Schemas and Database Row-Level Security}
\label{sec:rls}

\paragraph{Multi-schema mode.} With \code{db_schemas: [people, acad, fin, ...]} every table is named
\code{schema.table} in the prompt, the policy and the enforcer, and grants accept \code{fin.*}. A bare
name that exists in exactly one visible schema is completed automatically; one that exists in several
is returned to the model as a fixable error. Schema extraction reads PostgreSQL's catalogs, so the
connecting login needs no table privileges.

\paragraph{Row-level security as a second line of defence.} The enforcer's row filters and the
database's policies say the same thing twice. For each query NatQuery sends the user's role and identity
as transaction-local settings and then drops to a low-privilege database role:

\begin{lstlisting}[caption={What the database receives before the enforced query.}, label={lst:rls}]
SELECT set_config('erp.role', 'student', true);
SELECT set_config('erp.student_id', '1487', true);
SET LOCAL ROLE erp_ro;
-- enforced query runs here
\end{lstlisting}

A policy on each protected table reads these settings, so a bug in the application layer would still be
met by the database (Listing~\ref{lst:rlspol}). With no settings, a role sees no rows.

\begin{lstlisting}[caption={A row-level security policy of the ERP (excerpt).}, label={lst:rlspol}]
CREATE POLICY erp_ro_scope ON lib.loan FOR SELECT TO erp_ro
  USING (person_id = auth.ctx_int('person_id')
         OR auth.ctx_role() IN ('library', 'admin'));
\end{lstlisting}

Row-level security covers rows only. Column hiding and masking live in the enforcer, except for the
secret columns, which the database role is never granted.

\section{Execution Safeguards}

Table~\ref{tab:exec} lists what bounds the execution of a validated query. All values are defaults that
the policy can change.

\begin{table}[htbp]
\centering
\caption{Execution safeguards.}
\label{tab:exec}
\small
\renewcommand{\arraystretch}{1.15}
\begin{tabularx}{\textwidth}{L{4.2cm} Y}
\toprule
\hdr \textbf{Safeguard} & \textbf{Detail} \\
\midrule
Read-only transaction & Always rolled back, never committed \\
Statement timeout & 5 seconds by default (\code{statement_timeout_ms}) \\
Row cap & 500 rows by default, enforced in the SQL and again when fetching \\
Pooled connections & Connections are checked out per request and reset afterwards \\
Database role & Optional \code{SET LOCAL ROLE} to a SELECT-only role \\
Rate limits & 20 questions per minute and 120 model calls per hour per user by default \\
Generic user errors & Users see a generic message and a reference id; details go to the audit log; database errors are sanitised before reaching the model \\
\bottomrule
\end{tabularx}
\end{table}

\section{The Audit Log}

Every request writes JSON lines to an append-only file, each carrying the hash of the previous record.
With records $r_i$, genesis hash $h_0$ (64 zeros) and canonical JSON encoding,
\[
  h_i = \mathrm{SHA\text{-}256}\bigl(h_{i-1} \,\Vert\, r_i\bigr),
\]
or $\mathrm{HMAC}_K$ with SHA-256 when the key $K$ is set in the environment. Editing, deleting or
re-ordering any past record breaks the chain, which \code{natquery audit verify} checks. The
``authorized'' record is written \emph{before} the query runs, and if that write fails the query does
not run. The chain is tamper-\emph{evident}, not tamper-proof: without a key, someone who can rewrite
the whole file can forge a consistent chain, and cutting records off the end is not detectable by the
chain alone, so the log should be shipped off the machine.

\section{Authentication and Sessions}

Who the user is comes from an authentication provider behind a small interface; the built-in provider
keeps local accounts, and others (LDAP, single sign-on, the ERP's own session) can be registered.
Table~\ref{tab:auth} lists the controls of the built-in provider.

\begin{table}[htbp]
\centering
\caption{Authentication and session controls.}
\label{tab:auth}
\small
\renewcommand{\arraystretch}{1.15}
\begin{tabularx}{\textwidth}{L{4.2cm} Y}
\toprule
\hdr \textbf{Control} & \textbf{Detail} \\
\midrule
Password storage & scrypt (memory-hard, from the Python standard library), constant-time comparison \\
Account lockout & After repeated failures an account is locked for a period \\
No account enumeration & Wrong password and unknown user give the same message \\
Identity attributes & Roll number, employee id, department and so on are attached to the identity and bound into row filters \\
Sessions & Idle timeout (15 minutes) and absolute timeout (8 hours) by default; tokens are opaque; identity re-read on every request \\
Role changes & A role change or disabled account takes effect on the next request \\
\bottomrule
\end{tabularx}
\end{table}

\FloatBarrier
